% Zumdahl Chapter 7 free-response set -- 2 long (10) + 3 short (4) = 32 pts.
\documentclass{apchem-exam}

\apcunitword{Chapter}
\apcunit{7}
\apcunittitle{Atomic Structure and Periodicity}
\apcdoctitle{Free-Response Set}
\apcsubtitle{Zumdahl \S7.1--\S7.12 \textbullet\ 5 questions \textbullet\ 32 points \textbullet\ 50 minutes}

\begin{document}
\apctitleblock

\begin{apcnote}[Scope --- two CED exclusions apply]
Covers \textbf{CED 1.5} (\S7.5, \S7.8--7.11), \textbf{1.6}
(photoelectron spectroscopy, \S7.12), \textbf{1.7} (periodic trends,
\S7.12), \textbf{3.11} (\S7.1) and \textbf{3.12} (\S7.2).

Two exclusions shape what is asked. Topic 1.5 excludes
\textbf{assigning quantum numbers to individual electrons}, so \S7.6 is
not examined --- you will never be asked for $n$, $\ell$, $m_\ell$,
$m_s$. Topic 1.7 excludes \textbf{electron configurations of the aufbau
exceptions}, so \ce{Cr} and \ce{Cu} do not appear as configuration
questions. Both remain worth understanding; neither will be tested.

$h = \SI{6.626e-34}{\joule\second}$, \;
$c = \SI{3.00e8}{\metre\per\second}$.
\end{apcnote}

\examsection{II}{Free Response}{5 questions \textbullet\ 32 points \textbullet\ 50 minutes}{\calculatorok}

% =====================================================================
\frq{long}{10}{Zumdahl \S7.11--\S7.12 \textbullet\ CED 1.5, 1.6, 1.7 \textbullet\ configuration, PES and trends}

\begin{parts}
  \item Write the complete electron configuration of a phosphorus atom
        ($Z = 15$). \answerlines[2]{\ce{1s^2 2s^2 2p^6 3s^2 3p^3}}
  \item A photoelectron spectrum of phosphorus shows five peaks. State how
        many peaks correspond to each subshell and give the relative
        height of each. \answerlines[3]{One peak per occupied subshell:
        \ce{1s}, \ce{2s}, \ce{2p}, \ce{3s}, \ce{3p}. Relative heights
        follow the number of electrons: $2 : 2 : 6 : 2 : 3$ --- which sums
        to 15, the atomic number.}
  \item State which peak appears at the highest binding energy, and
        explain why. \answerlines[3]{The \ce{1s} peak. Those electrons are
        closest to the nucleus and are shielded by no other electrons, so
        they experience the largest effective nuclear charge at the
        smallest distance. The Coulombic attraction is therefore strongest
        and the most energy is needed to remove them.}
  \item The first ionization energy of sulfur is \emph{lower} than that of
        phosphorus, breaking the general left-to-right trend. Explain.
        \workspace[3cm]
        \keyonly{\begin{apcanswer}Phosphorus is \ce{3p^3}: each of the
        three \ce{3p} orbitals holds one electron. Sulfur is \ce{3p^4}, so
        one \ce{3p} orbital must hold a \emph{pair}.

        Two electrons confined to the same orbital repel each other
        strongly. That repulsion raises the energy of the paired electron
        and makes it easier to remove than the nuclear charge alone would
        suggest --- enough to more than offset the extra proton in sulfur.

        Note this is an argument about \textbf{electron--electron
        repulsion within one orbital}, not about ``half-filled shells
        being stable''.\end{apcanswer}}
\end{parts}

\begin{rubric}
  \rpoint{2}{(a) 2 pts fully correct; 1 pt if only the final superscript
    is wrong. Note the CED exclusion on 1.7: do not ask for \ce{Cr} or
    \ce{Cu}}
  \rpoint{2}{(b) 1 pt five peaks matched to the five subshells; 1 pt the
    $2\!:\!2\!:\!6\!:\!2\!:\!3$ heights}
  \rpoint{3}{(c) 1 pt \ce{1s}; 1 pt closest to the nucleus / least
    shielded; 1 pt greatest effective nuclear charge and therefore
    strongest attraction. An answer saying only ``it is the innermost''
    earns 1}
  \rpoint{3}{(d) 1 pt identifying the \ce{3p^3} versus \ce{3p^4}
    difference; 1 pt naming the pairing in one orbital; 1 pt the
    repulsion makes removal easier. ``Half-filled subshells are extra
    stable'' with no repulsion argument earns \textbf{0} for the final
    two points}
\end{rubric}

% =====================================================================
\frq{long}{10}{Zumdahl \S7.1--\S7.4 \textbullet\ CED 3.11, 3.12 \textbullet\ photons and atomic spectra}

Excited hydrogen atoms emit a bright blue-green line at
\SI{486}{\nano\metre}.

\begin{parts}
  \item Calculate the energy of a single photon of this light.
        \workspace[2.8cm]
        \keyonly{\begin{apcanswer}$E = \dfrac{hc}{\lambda} =
        \dfrac{(6.626\times10^{-34})(3.00\times10^{8})}
        {486\times10^{-9}} = \mathbf{4.09\times10^{-19}}$~J\end{apcanswer}}
  \item Calculate the energy emitted per mole of photons, in
        \si{\kilo\joule\per\mole}. \answerlines[2]{$(4.09\times10^{-19})
        (6.022\times10^{23}) = \num{2.46e5}$~J/mol
        $= \mathbf{246}$~kJ/mol}
  \item Hydrogen also emits a violet line at \SI{434}{\nano\metre}. State
        whether its photons carry more or less energy than those at
        \SI{486}{\nano\metre}, and justify without calculating.
        \answerlines[2]{\textbf{More.} Energy is inversely proportional to
        wavelength ($E = hc/\lambda$), and \SI{434}{\nano\metre} is the
        shorter wavelength.}
  \item Explain why the emission spectrum of hydrogen consists of discrete
        lines rather than a continuous band. \workspace[3cm]
        \keyonly{\begin{apcanswer}The electron in a hydrogen atom may
        occupy only certain allowed energy levels --- its energy is
        \textbf{quantized}. Emission occurs when the electron drops from
        one allowed level to another, and the photon carries away exactly
        the difference between them.

        Because only certain differences exist, only certain photon
        energies are possible, and therefore only certain wavelengths
        appear. A continuous spectrum would require every energy
        difference to be available, which would mean the electron could
        have any energy at all.\end{apcanswer}}
  \item State what the existence of these lines tells you about the energy
        of the electron in a hydrogen atom.
        \answerlines[2]{That it is restricted to discrete, quantized
        values rather than varying continuously.}
\end{parts}

\begin{rubric}
  \rpoint{3}{(a) 1 pt $E = hc/\lambda$; 1 pt wavelength converted to
    metres; 1 pt \SI{4.09e-19}{\joule}. Leaving $\lambda$ in nanometres
    gives an answer out by $10^{9}$ and earns at most 1}
  \rpoint{2}{(b) 1 pt multiplying by Avogadro's number; 1 pt
    \SI{246}{\kilo\joule\per\mole}}
  \rpoint{2}{(c) 1 pt ``more''; 1 pt the inverse relationship with
    wavelength}
  \rpoint{2}{(d) 1 pt energy levels are quantized; 1 pt the photon carries
    exactly the difference between two levels}
  \rpoint{1}{(e) the electron's energy is quantized}
\end{rubric}

% =====================================================================
\frq{short}{4}{Zumdahl \S7.12 \textbullet\ CED 1.7 \textbullet\ periodic trends}

\begin{parts}
  \item Arrange \ce{Na}, \ce{Mg} and \ce{K} in order of increasing atomic
        radius. \answerlines[2]{\ce{Mg} $<$ \ce{Na} $<$ \ce{K}}
  \item Justify your ordering using Coulombic reasoning.
        \workspace[3cm]
        \keyonly{\begin{apcanswer}\textbf{\ce{Mg} versus \ce{Na}:} same
        shell, but magnesium has one more proton. The extra nuclear charge
        pulls the same valence shell in more tightly, so \ce{Mg} is
        smaller.

        \textbf{\ce{K} versus \ce{Na}:} potassium's valence electron
        occupies the next shell out. It is both further from the nucleus
        and shielded by an additional filled shell, so the attraction it
        feels is weaker and the atom is larger.\end{apcanswer}}
\end{parts}

\begin{rubric}
  \rpoint{1}{(a) the order \ce{Mg} $<$ \ce{Na} $<$ \ce{K}}
  \rpoint{3}{(b) 1 pt increased nuclear charge across a period at constant
    shell; 1 pt a new shell further out down a group; 1 pt increased
    shielding. An answer resting on ``more electrons so bigger'' earns
    \textbf{0} --- it predicts \ce{Mg} $>$ \ce{Na}, which is wrong}
\end{rubric}

% =====================================================================
\frq{short}{4}{Zumdahl \S7.12 \textbullet\ CED 1.6 \textbullet\ reading a PES spectrum}

A photoelectron spectrum of an unknown element from period 3 shows peaks
at relative heights $2$, $2$, $6$ and $1$, from highest binding energy to
lowest.

\begin{parts}
  \item Identify the element, showing your reasoning.
        \answerlines[3]{The heights give the electron count in each
        subshell: $2 + 2 + 6 + 1 = 11$ electrons, so $Z = 11$ and the
        element is \textbf{sodium}. The configuration read off the
        spectrum is \ce{1s^2 2s^2 2p^6 3s^1}.}
  \item State which peak lies at the lowest binding energy, and what that
        indicates chemically. \answerlines[3]{The height-1 peak, the
        \ce{3s} electron. It is the outermost and most shielded electron,
        so it is held least tightly --- which is why sodium loses exactly
        one electron so readily and forms \ce{Na+}.}
\end{parts}

\begin{rubric}
  \rpoint{2}{(a) 1 pt summing the heights to 11 electrons; 1 pt sodium}
  \rpoint{2}{(b) 1 pt the \ce{3s} peak identified as lowest binding
    energy; 1 pt linking it to the loss of one electron / formation of
    \ce{Na+}}
\end{rubric}

% =====================================================================
\frq{short}{4}{Zumdahl \S7.2 \textbullet\ CED 3.12 \textbullet\ properties of photons}

\begin{parts}
  \item Light of \SI{300}{\nano\metre} strikes a metal whose electrons are
        bound with an energy of \SI{5.0e-19}{\joule}. Determine whether an
        electron is ejected. \workspace[2.8cm]
        \keyonly{\begin{apcanswer}$E_{\text{photon}} = \dfrac{hc}{\lambda}
        = \dfrac{(6.626\times10^{-34})(3.00\times10^{8})}
        {300\times10^{-9}} = \num{6.63e-19}$~J.

        This exceeds the \SI{5.0e-19}{\joule} binding energy, so an
        electron \textbf{is} ejected, carrying away the difference,
        about \SI{1.6e-19}{\joule}, as kinetic
        energy.\end{apcanswer}}
  \item Doubling the intensity of the light does not change whether an
        electron is ejected. Explain. \answerlines[2]{Intensity sets how
        \emph{many} photons arrive, not how much energy each one carries.
        Ejection depends on a single photon having enough energy on its
        own; more photons of the same energy eject more electrons, but
        cannot help if each is individually too weak.}
\end{parts}

\begin{rubric}
  \rpoint{2}{(a) 1 pt photon energy \SI{6.63e-19}{\joule}; 1 pt the
    comparison and conclusion that ejection occurs}
  \rpoint{2}{(b) 1 pt intensity means number of photons; 1 pt the
    one-photon-per-electron argument}
\end{rubric}

\examtotal

\end{document}
